%% P22 \dash{} Probability and Bayes' theorem
%%
%% Stub. The brief below is the contract for this program: what it must argue,
%% and what it must buy the reader. Writing the program means deleting the
%% \programstub{} block and replacing it with the program. If the brief turns
%% out to be wrong once the mathematics has been worked, change it and record
%% the contradiction in CLAUDE.md under *Resolved questions*.

\program{Probability and Bayes' theorem}
\label{prog:P22}

\programstub{%
Argument: probability is a measure on a sample space obeying three rules;
conditioning is restricting the space; Bayes' theorem is one line of algebra
from the definition and is derived, never asserted. Payoff: the base-rate
calculation an engineer must be able to do in a meeting \dash{} a classifier
with 99\% accuracy on a fault that occurs once in a thousand requests, and
what fraction of its alarms are real. Independence and conditional
independence distinguished, because the naive Bayes assumption and most
\enquote{these two evals are independent signals} claims live on the difference.
Trap: the prosecutor's fallacy, elicited from the reader before it is named.

\medskip
\textbf{Planned length:} 55 frames.
}
